\section{Introduction}
\label{sec:intro}

Modular blockchain architectures separate execution, settlement, and data
availability into distinct layers, each maintained by a different network. A
rollup or sovereign chain typically settles finality claims on a base layer
such as Bitcoin and publishes transaction data to a DA network such as
Celestia~\cite{albassam2019lazyledger}. Suppose such a chain makes block
validity depend on those layers: a block is valid only if it carries a fresh
settlement checkpoint and a DA inclusion receipt. The chain then inherits
their liveness. If Bitcoin settlement is congested or censored, or the DA
network is unavailable, no proposal can satisfy the validity predicate. The
BFT consensus layer~\cite{buchman2016tendermint,buchman2018gossip} built on
top stalls even though the validator set is fully connected and correct.
Dropping the dependency avoids the stall but gives up the settlement and
availability guarantees that motivated the modular design.

\textbf{The gap.} Existing designs resolve this tension outside the
validator set's own agreement. Some ignore health and assume the periphery is
always available. Others handle it out of band, with operators intervening
manually. L2 rollups delegate it to an asynchronous fallback path on another
chain, which users trigger only after the primary path has stalled (escape
hatches and forced inclusion, Section~\ref{sec:related}). Protocols that
trade finality for availability inside
consensus~\cite{neu2021ebbflow,buterin2020gasper} respond to \emph{internal}
participation, not to the health of \emph{external} infrastructure. To our
knowledge, no design lets the validator quorum agree on external-layer health
in the same vote that orders transactions, and then change the chain's
validity rules and withdrawal policy on the basis of that agreement.

\textbf{Why the decision belongs in consensus.} The circuit breaker and the
relaxation of receipt requirements change two things that every validator
must agree on: whether a block is valid, and whether a withdrawal executes.
Suppose each validator applied a local policy to its own peripheral view.
Validators with different views would then disagree on validity, so proposals
would collect nil votes and rounds would repeat, or they would execute
withdrawals differently and produce diverging \texttt{AppHash} values.
Engram therefore reads sensors only inside vote extensions and writes FSM
state only from a certificate that a quorum has already agreed on
(Section~\ref{subsec:health-certificate}). An escape hatch does not provide
this property: it runs on another chain, is triggered by users or watchers
after the chain has stalled, and does not keep the chain live.

\textbf{What Engram is -- and is not.} Engram is not a new BFT agreement
algorithm. Proposer selection, the prevote/precommit protocol, quorum
intersection, and deterministic finality are provided unchanged by a
CometBFT core, and Engram inherits their safety and liveness conditions.
Engram's contribution sits one level above and extends three things:
(i) the \emph{consensus object}, with a quorum-certified health certificate
and an FSM state; (ii) the \emph{proposal-validation predicate}, so that a
proposal is valid only if that state follows deterministically from the
certificate; and (iii) the \emph{state-transition logic}, so that the chain's
validity requirements (Bitcoin and DA receipts) and its withdrawal policy
change with that state. We use ``Hybrid Adaptive Consensus'' in this sense: an adaptive
layer integrated into the consensus pipeline, not a replacement for it.

\textbf{Scope of ``degradation''.} We distinguish two kinds of network.
The \emph{core consensus network} carries CometBFT's own votes. Like any BFT
protocol, Engram requires that at least $2f{+}1$ validators can communicate
after global stabilization time; no state machine can restore liveness if
that fails. The \emph{peripheral networks} are Bitcoin,
Celestia, and the P2P interfaces validators use to reach them, together with
peer-quality signals (churn, latency, subnet diversity) on the Engram P2P
interface that fall short of breaking core quorum connectivity. Engram
handles degradation of the latter while the former holds. Throughout the
paper, ``P2P degradation'' refers to this peripheral or partial degradation,
never to a partition of the core quorum.

\textbf{Contributions.}
\begin{itemize}
    \item \textbf{Quorum-certified peripheral health.} A health certificate
    built from signed, height-bound observations of at least $2f{+}1$ voting
    power, carried in CometBFT vote extensions. A deterministic
    weighted-median/majority function aggregates them, and its output is
    bounded by honest observations whichever observations a Byzantine
    proposer includes or omits
    (Section~\ref{subsec:health-certificate}).
    \item \textbf{Adaptive fallback state machine.} A four-state FSM
    (\textsc{Anchored}, \textsc{Suspicious}, \textsc{Sovereign},
    \textsc{Recovering}) evaluated on the certificate, with asymmetric
    hysteresis, exponential backoff on repeated failed recovery, and an
    economic circuit breaker~\cite{mendonca2020resiliency} that locks
    cross-chain withdrawals while the chain runs without external
    corroboration (Sections~\ref{subsec:fsm}--\ref{subsec:transitions}).
    \item \textbf{Validator-verified recovery proof.} A Noir/UltraHonk circuit
    that proves header-chain continuity, FSM-transition legality,
    withdrawal-lock adherence, and boundary-validator signatures for up to
    256 \textsc{Sovereign}-interval headers at constant proving and
    verification cost. Longer intervals use a sequence of proofs linked
    through an on-chain committed root, optionally folded by an aggregator
    (Appendix~\ref{app:aggregation}). Validators verify each proof inside
    block execution; Bitcoin timestamps the resulting checkpoint but does not
    verify the proof (Section~\ref{subsec:reanchoring}).
    \item \textbf{Layered formal specification.} A four-layer TLA+
    specification. Bounded model checking of each layer's safety and liveness
    properties found zero violations at $N{\le}4$; the largest single run
    explored 3.39 million distinct states. In a five-mechanism ablation
    study, removing each mechanism reproduces a distinct predicted
    counterexample. The cross-layer refinement obligation is specified but
    open (Section~\ref{sec:formal}).
    \item \textbf{Live evaluation.} A 7-validator Cosmos SDK testnet against
    live Bitcoin regtest and Celestia networks with network- and
    application-level fault injection, comparing Engram against a
    hard-gated modular baseline and a vanilla CometBFT control
    (Sections~\ref{sec:setup}--\ref{sec:evaluation}).
\end{itemize}

\begin{table}[h]
\caption{Research Questions}
\label{tab:rq}
\centering
\small
\begin{tabular}{@{}p{0.07\linewidth}p{0.85\linewidth}@{}}
\toprule
\textbf{RQ} & \textbf{Question} \\
\midrule
RQ1 & \textbf{Safety.} Under up to $f$ colluding Byzantine validators, does
certifying peripheral health by quorum prevent conflicting FSM states,
forged or stale receipts, false health claims, and withdrawal leakage? \\
RQ2 & \textbf{Liveness.} When Bitcoin or DA becomes unavailable, or the
peripheral P2P view degrades while core quorum connectivity holds, does
Engram keep committing blocks where a hard-gated modular baseline stalls? \\
RQ3 & \textbf{Recovery.} Once peripherals recover, do hysteresis and the
recovery proof return the chain to a stable \textsc{Anchored} state without
oscillation? \\
RQ4 & \textbf{Cost.} What overhead do the extended consensus object,
certificate validation, and recovery proof add, and where does it come
from: proposal size, validation time, or additional consensus rounds? \\
\bottomrule
\end{tabular}
\end{table}

The remainder of this paper is organized as follows. Section~\ref{sec:related}
situates Engram in the literature. Section~\ref{sec:design} presents the
protocol, including the health certificate and the recovery proof.
Section~\ref{sec:formal} presents the formal specification and model-checking
results. Section~\ref{sec:setup} describes the testbed and baselines,
Section~\ref{sec:evaluation} answers RQ1--RQ4, Section~\ref{sec:discussion}
discusses limitations, and Section~\ref{sec:conclusion} concludes.
